This chapter introduces the theoretical machinery used in the rest
of the thesis: the structured-prediction setting and the loss
functions (Section~\ref{sec:structured-prediction}); Markov Networks as the
structured decision layer (Section~\ref{sec:markov-networks}); MAP
inference on chains and on cyclic graphs
(Section~\ref{sec:map-inference}); max-margin learning with the
LP-relaxed surrogate (Section~\ref{sec:learning-mn}); and the neural
feature extractors that supply the unary potentials
(Section~\ref{sec:nn-features}). The notation used throughout is
summarized below.

\begin{center}
\renewcommand{\arraystretch}{1.15}
\begin{tabular}{l l}
\hline
\textbf{Symbol} & \textbf{Meaning} \\
\hline
$\mathcal{X},\ \mathcal{Y}$ & input and output (label) sets \\
$K = |\mathcal{Y}|$ & number of labels \\
$N = |V|$ & number of output variables \\
$T$ & chain length (chain graphs only; $T = N$ for chains) \\
$x,\ y$ & input and joint output configuration \\
$y_v$ & label at node $v$ \\
$G = (V, E)$ & graph of the Markov network \\
$\mathcal{N}(v)$ & set of neighbours of node $v$ in $G$ \\
$f(x, y)$ & joint score (unnormalized log-probability) \\
$f_v(x, y_v)$ & unary potential at node $v$ \\
$f_{vv'}(y_v, y_{v'})$ & pairwise potential at edge $\{v, v'\}$ \\
$\psi(x, y)$ & joint feature map \\
$w \in \mathbb{R}^d$ & Markov-network parameter vector \\
$\theta$ & neural backbone parameters \\
$g_\theta$ & neural backbone (unary feature extractor) \\
$W \in \mathbb{R}^{K \times K}$ & learned pairwise potential matrix (shared across edges) \\
$\mu_v(a),\ \mu_{vv'}(a, b)$ & relaxed LP indicators \\
$\phi^{i}$ & per-example dual variables in LP-M3N \\
$\Delta,\ \Delta^{\mathrm{LP}}$ & margin-rescaling loss and LP-relaxed surrogate \\
$\ell$ & task loss used at training (Hamming) \\
$L_{0/1},\ L_H$ & evaluation losses (0/1, Hamming) \\
$\lambda$ & regularization strength \\
$m$ & training-set size \\
$\mathcal{D}$ & training set $\{(x^i, y^i)\}_{i=1}^{m}$ \\
$\hat{y}$ & MAP-predicted labeling \\
\hline
\end{tabular}
\end{center}

\bigskip

\section{Structured Prediction}
\label{sec:structured-prediction}

Standard classification maps an input $x$ to a single label from a
finite set. In many real-world problems, however, the output is not a
single label but a composite object whose components are interdependent.
Parsing a sentence yields a tree, segmenting an image yields a labeling
of every pixel, and solving a Sudoku puzzle yields a $9 \times 9$ grid
in which every cell label is constrained by the labels of every other
cell in its row, column, and box. Problems of this form are known as
\emph{structured output prediction}~\cite{Bakir2007}.

\subsection{Formal setting}

Let $\mathcal{X}$ denote the input space, $\mathcal{Y} = \{1, \ldots,
K\}$ a finite per-component label set, and $N$ the number of output
components. A structured-output predictor is a function
$h\colon \mathcal{X} \to \mathcal{Y}^{N}$ that maps each input to a
joint configuration $y = (y_1, \ldots, y_N)$. Two features
distinguish the structured case from ordinary classification:
\begin{itemize}
    \item The joint output space $\mathcal{Y}^{N}$ grows exponentially
    with the number of output variables. For a Sudoku grid with
    $N = 81$ cells and $K = 9$ possible digits,
    $|\mathcal{Y}^{N}| = 9^{81} \approx 2 \cdot 10^{77}$ raw
    configurations, even before any constraint is imposed. Enumerating
    $\mathcal{Y}^{N}$ is therefore infeasible.
    \item Output components are not independent. The quality of a
    configuration is a joint quantity that couples several variables at
    once.
\end{itemize}

A convenient way to handle both properties is to introduce a scoring
function $f\colon \mathcal{X} \times \mathcal{Y}^{N} \to \mathbb{R}$
and to define the predictor as
\begin{equation}
    h(x) \;=\; \argmax_{y \in \mathcal{Y}^{N}} f(x, y).
    \label{eq:map-prediction}
\end{equation}

Equation~\eqref{eq:map-prediction} is referred to as \emph{MAP
prediction} (maximum-a-posteriori), since in probabilistic models
$f(x, y)$ plays the role of an unnormalized log-probability. For this
formulation to be useful in practice, the scoring function $f$ must
admit a compact parameterization that captures the dependencies among
the output components, and the argmax must be solvable without
enumerating $\mathcal{Y}^{N}$ (Section~\ref{sec:map-inference}).


\subsection{Loss functions}
\label{subsec:loss-functions}

Learning and evaluation both require a notion of how wrong a prediction
$\hat{y}$ is relative to a ground-truth output $y$. The structured
problems considered in this thesis decompose into $N$ components,
$y = (y_1, \ldots, y_N)$ with $y_i \in \{1, \ldots, K\}$. Two loss
functions are used throughout.

\paragraph{0/1 loss.}
The strict all-or-nothing loss counts a prediction as correct only when
every component matches:
\begin{equation}
    L_{0/1}(y, \hat{y}) \;=\; \mathbb{I}\!\left[y \neq \hat{y}\right]
    \;=\; 1 - \prod_{i=1}^{N} \mathbb{I}\!\left[y_i = \hat{y}_i\right],
    \label{eq:zero-one-loss}
\end{equation}
which equals $1$ whenever the predicted vector differs from the
ground truth in at least one position, and $0$ otherwise. This is
the natural measure when a partially correct answer is useless: a
Sudoku grid with a single wrong digit is not a valid solution.

\paragraph{Hamming loss.}
The Hamming loss counts the number of component-wise mismatches,
\begin{equation}
    L_{H}(y, \hat{y}) \;=\; \sum_{i=1}^{N}
    \mathbb{I}\!\left[y_i \neq \hat{y}_i\right],
    \label{eq:hamming-loss}
\end{equation}
and the \emph{normalized} Hamming loss divides by $N$ to obtain a
quantity in $[0, 1]$:
\begin{equation}
    L_{H}^{\mathrm{norm}}(y, \hat{y}) \;=\;
    \frac{1}{N} \sum_{i=1}^{N}
    \mathbb{I}\!\left[y_i \neq \hat{y}_i\right].
    \label{eq:hamming-loss-norm}
\end{equation}
Hamming loss is a finer-grained measure that rewards partial
correctness. Together with the 0/1 loss it forms the evaluation pair
used in Chapter~\ref{chap:experiments}.

Throughout this thesis, the normalized Hamming loss serves a dual
role: as the task loss $\ell$ used inside max-margin training
(Section~\ref{sec:learning-mn}), and as one of the evaluation losses
reported in Chapter~\ref{chap:experiments}. The 0/1 loss is used at
evaluation time only --- it is not node-decomposable and is
incompatible with margin-rescaling at training time.

\subsection{Risk minimization}

The goal of learning is to select a predictor $h$ that minimizes the
\emph{expected risk}
\begin{equation}
    R(h) \;=\; \mathbb{E}_{(x, y) \sim p}\!\left[
    L\bigl(y, h(x)\bigr) \right],
    \label{eq:expected-risk}
\end{equation}
where $p$ is the unknown data-generating distribution and $L$ is a
task-specific loss. Since $p$ is not available, one works with a
training sample
$\mathcal{D} = \{(x^i, y^i)\}_{i=1}^{m}$ drawn i.i.d.\ from $p$ and
minimizes the \emph{empirical risk}
\begin{equation}
    \hat{R}(h) \;=\; \frac{1}{m} \sum_{i=1}^{m}
    L\bigl(y^i, h(x^i)\bigr).
    \label{eq:empirical-risk}
\end{equation}

In practice, the labeled data are split into three disjoint
subsets: a \emph{training set} on which the empirical risk (or its
surrogate) is minimized, a \emph{validation set} used to select
hyperparameters such as the regularization strength of
Section~\ref{sec:learning-mn} and to monitor overfitting, and a
\emph{test set}, held out until the end, that yields an unbiased
estimate of the expected risk.

Directly minimizing the 0/1 or Hamming loss over a structured output
space is generally intractable: the losses are discrete, and the
hypothesis class defined by Eq.~\eqref{eq:map-prediction} is piecewise
constant in its parameters. Practical learning algorithms therefore
minimize a convex \emph{surrogate} loss --- such as the max-margin loss
discussed in Section~\ref{sec:learning-mn} --- while aiming to control
the original risk.
